\subsection{判别器留出诊断与自建程序行为试验}
为检查新判别器，构造60个父配置，每个含10种版本，共600例；
按父配置固定40个开发、20个留出，留出200例含七类已知故障各20例、40个无害对照和20个不支持解析器例。
故障为未知参数、缺必填、非法枚举、错误类型、缺入口、资源变化和语法错误。
这些是研究者控制的机械干预，不是60个天然社区技能，也不是由独立专家给出的整体质量金标准。
对照含仅参数名、标准库argparse、文档均分阈值及去模块消融。
文档均分阈值0.8为本研究的工程适配，不能作为发表版SkillNet评价器的复现。

行为试验使用整数求和、稳定全行去重及DAG判定三个明确任务族，每族四种自建实现，共12个程序；
每个实现固定24个输入，共288次真实输出。
三种实现满足登记套件，九种实现植入可观测缺陷；所有文档都可取得约0.92的同类启发评分。
执行代码和输出回执均有字节指纹，验收器与被测实现分离。
这些观察验证“有相同高文档分的不同实现仍可由产物区分”，不估计自然技能优劣的发生率。

\subsection{公开数据派生的18节点可执行网络}
为从历史检索覆盖推进到真正路径执行，另建立一个可审计的18节点实验网络：
解析、按键去重、跨分割排除、空分数删除、指标计算及三阶段复合清洗各三个版本。
节点包括正确、条件受限和未探测实现；公开说明对同能力保持中性，ID为不透明散列。
选择接口只接收ID、能力、文本、资源指纹及字节成本，拒收测试答案、输出、实现键和缺陷标签。
不执行社区包；工作进程只允许本研究固定编写的算子。

原始表来自UCI：开发使用Iris（150行）和Wine（178行），留出使用WDBC（569行）与Banknote（1,372行）
\citep{UCIiris,UCIwine,UCIwdbc,UCIbanknote}。
各原始文件、下载归档、DOI、CC BY 4.0许可说明及指纹随稿保存。
将各文件第一数值特征归一化为四位小数分数、目标映射为二元标签，
再显式加入冲突重复、跨分割主体、空字段、含逗号及换行的引号字段和0.5阈值边界。
它们是\emph{公开数据派生的构造审计任务}，不代表原数据集分类准确率，不涉及临床判断。
八种清洗需求、24/72两个行数、两个引号条件及两个边界条件，形成每数据集64个任务。
开发128个、主留出128个，另从留出材料构造24个高精度数值压力输入。
任务语义和算子库跨数据集共享，故这种划分不是跨任务族或自然技能来源的独立验证。

独立参照按索引与自写CSV状态机读取，指标采用精确有理数；
worker使用标准库CSV、Decimal或各受测变体。
参照不导入规划器或worker实现；规划器也不导入答案计算和算子函数。
两路仍共享研究者规定的任务语义，降低实现耦合不等于第三方真值。
开发集共实际执行928次单节点探测，测试结果不可用于这些证据。
方案、输入、代码快照、公开元数据及基线在测试执行前冻结；这不是外部时间戳预注册，
研究者知道构造规则，不能称为独立双盲测试。

\subsection{预算、主估计量与压力检验}
预算为每任务元数据最低可行路径字节的1.0、1.5及2.0倍，向上取整，至多五次算子调用。
不使用测试成功来挑预算。
主估计量先在每个留出任务内平均三档二元成功，再平均128个任务；两个数据集权重相同。
32个清洗--引号--阈值条件块各有两个来源、两个行数，按块进行10,000次配对条件重采样。
这只描述所选构造条件的变动；不作真实任务总体置信度、多重比较显著性或分布无关保证。

固定压力包括新数值精度、取证后节点换版本以及缺失全部必要指标节点。
另做三次不透明ID置换，仅检验命名/并列顺序敏感性，复用相同输入。
主运行共7,856个策略决定，完全相同输入及实现路径共享一次执行，实际为1,204次路径执行、3,934次算子调用。
因此7,856不是独立任务数，1,204也不是全部不同输入。
在主结果产生后，再将精度与版本变化合并作明确标注的事后证伪：24个已有精度输入、240个决定、132次路径执行；
没有用该反例修改方法或计入预定主终点。
该前一阶段自建判别器与18节点网络实验为零模型调用；本轮外部付费实验另行记录。
本机只在去凭据子进程运行自建代码；Linux CI另在无网络、只读、非root、去能力及限资源容器复现同批输入。
容器复现不是新增科学样本，也不是对社区恶意代码的安全认证。
